\Needspace{6\baselineskip}
\begin{center}\hypertarget{native-259-article186}{}\pdfbookmark[2]{186.}{article-186}186.\end{center}
Moreover, the principle that the squares of the differences between the observed and calculated quantities should yield the smallest possible sum can also be considered in the following way, independently of the calculus of probability.

Whenever the number of unknowns equals the number of independent observed quantities, we can determine the unknowns so as to satisfy the observations exactly. But when the former number is smaller than the latter, absolutely exact agreement cannot be obtained insofar as the observations lack absolute precision. In this case we must therefore strive for the best possible agreement, or to make the differences as small as possible. There is, however, something inherently vague about this notion. A system of values for the unknowns that makes \emph{all} the respective differences smaller than another is undoubtedly to be preferred to it. Even so, the choice between two systems, one of which gives better agreement for some observations and the other for others, is in a sense left to our discretion. Clearly, innumerable different principles can be proposed that satisfy the former condition. Denote the differences between observation and calculation by $\Delta$, $\Delta'$, $\Delta''$, etc. The former condition is satisfied when $\Delta\Delta+\Delta'\Delta'+\Delta''\Delta''+\text{etc.}$ is a minimum (which is our principle). It is also satisfied when $\Delta^4+\Delta'^4+\Delta''^4+\text{etc.}$ or $\Delta^6+\Delta'^6+\Delta''^6+\text{etc.}$ or, more generally, the sum of powers with any even exponent attains a minimum. Of all these principles, however, ours is the simplest, since the others would lead us to extremely complicated calculations. Moreover, our principle, which we have used since 1795, has recently also been put forward by the distinguished \textsc{Legendre} in his work \emph{Nouvelles méthodes pour la
\SourceLeaf{260}détermination des orbites des comètes, Paris} 1806. Several other properties of this principle are set out there, which we omit here for brevity.

If we adopted a power with an infinitely large even exponent, we would be led to the system in which the largest differences are made as small as possible.

To solve linear equations whose number exceeds the number of unknown quantities, the celebrated \textsc{Laplace} uses another principle that the distinguished \textsc{Boscovich} had already proposed, whereby the differences themselves, all taken with positive sign, are to yield the smallest sum. It is easy to show that a system of values for the unknowns derived from this principle alone must necessarily\footnote{Except in special cases where the solution remains in some measure undetermined.} satisfy exactly as many of the given equations as there are unknowns. The remaining equations thus enter into consideration only insofar as they \emph{help decide the choice}. Thus if, for example, the equation $V=M$ is among those that are not satisfied, the system of values found by that principle would remain unchanged even if, instead of $M$, any other value $N$ had been observed, provided only that, with $n$ denoting the calculated value, the differences $M-n$ and $N-n$ have the same sign. The celebrated \textsc{Laplace}, however, modifies this principle somewhat by adding a new condition. He requires that the sum of the differences themselves, with their signs unaltered, should become $=0$. As a result, the number of equations exactly satisfied becomes one less than the number of unknown quantities. Even so, our previous observation still holds, provided there are at least two unknowns.
